\chapter{Marchés et microstructure}\label{ch:iv2:markets-and-microstructure}

\begin{ivtrains}
Lire un carnet d'ordres, ce que rapporte la fourchette, et l'effet de ses propres ordres sur le cours~: files, sélection adverse, impact, TWAP, VWAP et liquidation optimale, vente à découvert, ainsi que les règles qui encadrent l'ensemble. Chaque énoncé donne ses chiffres. Les règles citées sont américaines, sauf si un énoncé précise le contraire.
\end{ivtrains}

\section{Lire le carnet d'ordres}

\begin{question}[$\star$ \iqroles{QT, Trader, QR, Dev} \iqfirm{teneur de marché, hedge fund} \iqlevel{junior}]\label{iq:iv2:markets-and-microstructure:1}
Un carnet d'ordres affiche des cours acheteurs de 300 actions à 99.98 et de 500 à 99.97, et des cours vendeurs de 100 actions à 100.02, de 400 à 100.03 et de 1\,000 à 100.05. Vous envoyez un ordre au marché d'achat de 300 actions. Quel est votre prix moyen, et combien payez-vous au-dessus du prix milieu, par action~? Que deviennent le meilleur cours vendeur et le prix milieu du carnet~?
\end{question}

\begin{question}[$\star$ \iqroles{QT, Trader, QR} \iqfirm{teneur de marché, hedge fund} \iqlevel{junior}]\label{iq:iv2:markets-and-microstructure:2}
Le meilleur cours acheteur est de 100.00 pour 900 actions, et le meilleur cours vendeur de 100.02 pour 100 actions. Quel est le prix milieu, et quel est le microprix pondéré par la taille, qui pondère le prix de chaque côté par la taille de l'autre côté~? Dans quel sens attendriez-vous le mouvement suivant, et pourquoi~?
\end{question}

\begin{question}[$\star$ \iqroles{QT, Dev, QD, Trader} \iqfirm{teneur de marché, hedge fund} \iqlevel{junior}]\label{iq:iv2:markets-and-microstructure:3}
À un même prix, trois ordres de vente en attente arrivent dans cet ordre~: $A$ pour 100 actions, $B$ pour 200 et $C$ pour 50. Un ordre d'achat de 250 actions arrive. Comment est-il réparti selon la priorité prix-temps~? Comment l'est-il dans le pur prorata, où chaque ordre est exécuté au prorata de sa taille~? Quelle règle récompense la rapidité, et laquelle récompense la taille~?
\end{question}

\begin{question}[$\star$ \iqroles{QT, Trader, QR} \iqfirm{teneur de marché, hedge fund, banque} \iqlevel{junior}]\label{iq:iv2:markets-and-microstructure:4}
Vous achetez à 100.03 alors que le prix milieu est de 100.00. Cinq minutes plus tard, le prix milieu vaut 100.02. Définissez et calculez, en cents, le spread effectif, le spread réalisé et l'impact sur le prix. Quelle part le teneur de marché a-t-il conservée~?
\end{question}

\begin{question}[$\star$ \iqroles{QT, Trader} \iqfirm{teneur de marché, hedge fund} \iqlevel{junior}]\label{iq:iv2:markets-and-microstructure:5}
Un titre a une fourchette de 1 cent. La plateforme facture au preneur de liquidité 0.3 cent par action et verse 0.2 cent au fournisseur de liquidité. Si vous voulez acheter, combien économisez-vous par action en vous faisant exécuter passivement au cours acheteur plutôt qu'en traversant la fourchette au cours vendeur~? Combien un teneur de marché gagne-t-il net par action, sur un aller-retour où les deux côtés sont exécutés passivement, sans sélection adverse~?
\end{question}

\begin{question}[$\star$ \iqroles{QT, Trader, QR} \iqfirm{teneur de marché, hedge fund} \iqlevel{junior}]\label{iq:iv2:markets-and-microstructure:6}
Avec un pas de cotation de 1 cent, un titre s'échange à 5~\$ et un autre à 500~\$. Quelle est la fourchette minimale de chacun, en points de base~? Lequel a le plus de chances de rester à une fourchette d'un tick avec une file longue, et pourquoi~? Pour le négocier, qu'est-ce qui compte le plus, la rapidité ou le prix~?
\end{question}

\begin{question}[$\star$ \iqroles{QR, QT} \iqfirm{teneur de marché, hedge fund} \iqlevel{junior}]\label{iq:iv2:markets-and-microstructure:7}
Les variations successives du prix d'un titre ont une autocovariance de $-0.0004$ dollars au carré à un décalage. Supposez que le prix efficient suit une marche aléatoire et que chaque transaction a autant de chances d'être au cours acheteur qu'au cours vendeur, indépendamment des autres. Quelle fourchette cela implique-t-il~? Pourquoi l'autocovariance est-elle négative du tout~?
\end{question}

\begin{question}[$\star$ \iqroles{QT, Trader, Analyste} \iqfirm{hedge fund, banque} \iqlevel{junior}]\label{iq:iv2:markets-and-microstructure:8}
Les transactions d'une journée sur un titre sont de 100 actions à 10.00, 200 à 10.20 et 700 à 10.10. Quel est le prix moyen pondéré par les volumes (VWAP)~? Vous avez acheté 500 actions dans la journée à 10.08 en moyenne. Quel a été votre résultat par rapport à lui, en points de base~? Que donne la moyenne simple des prix~?
\end{question}

\begin{question}[$\star$ \iqroles{QT, Trader, Analyste} \iqfirm{hedge fund, banque} \iqlevel{junior}]\label{iq:iv2:markets-and-microstructure:9}
Un titre compte 30 millions d'actions vendues à découvert, et 2 millions d'actions s'échangent chaque jour. Quels sont ses jours de couverture~? Si les vendeurs à découvert ne peuvent représenter que 20~\% du volume lorsqu'ils rachètent, combien de temps faut-il~? Pourquoi ce nombre compte-t-il lorsqu'un titre bondit~?
\end{question}

\begin{question}[$\star$ \iqroles{QT, Trader} \iqfirm{teneur de marché, hedge fund, banque} \iqlevel{junior}]\label{iq:iv2:markets-and-microstructure:10}
Un certificat de dépôt américain, titre négocié aux États-Unis, représente 2 actions ordinaires cotées à Londres. Les actions s'échangent à 10.00~\textsterling, et une livre vaut 1.25~\$. Quelle est la fair value du certificat~? Il s'échange à 25.20~\$, et l'ensemble des coûts de la transaction s'élève à 30 points de base. Que faites-vous, et combien gagnez-vous~?
\end{question}

\begin{question}[$\star$ \iqroles{QT, Dev, QD, Trader} \iqfirm{teneur de marché, hedge fund} \iqlevel{junior}]\label{iq:iv2:markets-and-microstructure:11}
Trois plateformes américaines affichent des cours vendeurs de 100 actions à 100.01 sur $A$, de 500 à 100.02 sur $B$ et de 400 à 100.03 sur $C$. Vous voulez 300 actions et vous privilégiez la plateforme $B$. Pouvez-vous simplement envoyer un ordre d'achat à $B$ à 100.02~? Qu'est-ce qu'un trade-through, et qu'est-ce qu'un ordre de balayage intermarchés (ISO)~?
\end{question}

\begin{question}[$\star$ \iqroles{QT, Trader, Dev} \iqfirm{teneur de marché, hedge fund} \iqlevel{junior}]\label{iq:iv2:markets-and-microstructure:12}
Une enchère d'ouverture comporte des ordres d'achat à cours limité~: 100 actions à 101, 200 à 100 et 300 à 99, et des ordres de vente à cours limité~: 150 à 99, 150 à 100 et 200 à 101. Quel prix maximise le volume échangé, et quel volume est échangé~? Qui ne s'exécute pas~?
\end{question}

\begin{question}[$\star$ \iqroles{QT, Trader} \iqfirm{hedge fund, banque} \iqlevel{junior}]\label{iq:iv2:markets-and-microstructure:13}
Vous devez acheter 200\,000 actions à 10~\% du volume total, vos propres transactions comprises. Les autres intervenants échangeront 2 millions d'actions aujourd'hui. Combien pouvez-vous faire aujourd'hui, et quelle part du volume des autres cela représente-t-il~? Qu'est-ce qui fait qu'un ordre au pourcentage du volume poursuit les cours~?
\end{question}

\begin{question}[$\star$ \iqroles{QT, Trader, Analyste, Risque} \iqfirm{hedge fund, banque} \iqlevel{junior}]\label{iq:iv2:markets-and-microstructure:14}
Vous voulez vendre à découvert 10\,000 actions d'un titre américain aujourd'hui. Que doit avoir fait votre courtier avant d'accepter l'ordre~? Qu'est-ce qu'un «~fail to deliver~», et une vente à découvert sans locate est-elle toujours illégale~? Citez une catégorie d'intervenants qui bénéficie d'une exception.
\end{question}

\begin{question}[$\star$ \iqroles{QT, QR, Trader} \iqfirm{hedge fund, teneur de marché} \iqlevel{junior}]\label{iq:iv2:markets-and-microstructure:15}
Prenez cette règle~: le coût moyen d'achat de $Q$ actions, avec un volume quotidien $V$ et une volatilité quotidienne $\sigma$, est d'environ $\sigma\sqrt{Q/V}$. Un titre a une volatilité quotidienne de 2~\%. Combien coûte l'achat de 1~\% du volume d'une journée, et de 4~\%, en points de base~? Si vous quadruplez votre ordre, que devient le coût en dollars~?
\end{question}

\section{Coûts, files et indices de référence}

\begin{question}[$\star\star$ \iqroles{QT, Trader, Analyste} \iqfirm{hedge fund, banque} \iqlevel{junior}]\label{iq:iv2:markets-and-microstructure:16}
Vous décidez d'acheter 10\,000 actions alors que le cours est de 50~\$. Vous n'en achetez que 8\,000, à 50.10~\$ en moyenne, et la commission est de 1 cent par action. Le cours clôture à 50.30~\$, et vous n'achetez jamais les 2\,000 autres. Quelle est l'implementation shortfall, en dollars et en points de base du portefeuille théorique~? Décomposez-la en ses éléments.
\end{question}

\begin{question}[$\star\star$ \iqroles{QT, QR, Dev} \iqfirm{teneur de marché, hedge fund} \iqlevel{junior, confirmé}]\label{iq:iv2:markets-and-microstructure:17}
Votre ordre rejoint une file où 20 ordres d'un seul lot le précèdent. Les événements qui retirent un ordre en tête de file, transactions ou annulations, arrivent selon un processus de Poisson à un taux de 1.5 par seconde. Considérez que les retraits ne proviennent que de la tête de file. Vous êtes exécuté par le retrait qui emporte votre propre ordre, c'est-à-dire celui qui survient après le départ des 20 ordres qui vous précèdent. Combien de temps, en moyenne, faut-il avant votre exécution~? Quel est l'écart-type~? Quelle est la probabilité d'être exécuté dans les 20 secondes~?
\end{question}

\begin{question}[$\star\star$ \iqroles{QT, QR} \iqfirm{teneur de marché, hedge fund} \iqlevel{junior, confirmé}]\label{iq:iv2:markets-and-microstructure:18}
Un titre a une fourchette de 1 cent, et vous postez au cours acheteur. Si vous êtes en tête de file, de petites transactions vous exécutent, et le cours baisse alors d'un tick dans 20~\% des cas. Si vous êtes en fin de file, vous n'êtes exécuté que lorsqu'un gros ordre arrive, et le cours baisse alors d'un tick dans 70~\% des cas. S'il ne baisse pas, il reste à son niveau. Quel gain par action espérez-vous, par rapport au prix milieu juste avant votre exécution, dans chacun des deux cas~? Quelle est la valeur de la position dans la file~?
\end{question}

\begin{question}[$\star\star$ \iqroles{QT, QR, Trader} \iqfirm{teneur de marché, hedge fund} \iqlevel{junior, confirmé}]\label{iq:iv2:markets-and-microstructure:19}
Un titre vaut soit 100, soit 120, avec 30~\% de chances de valoir 120. Un quart des traders connaissent la vraie valeur et échangent en conséquence. Les autres achètent ou vendent avec la même probabilité. Chacun échange une action, et la concurrence vous amène à l'équilibre sur chaque cotation. Quels sont vos cours acheteur et vendeur~? Pourquoi ne sont-ils pas symétriques autour de la moyenne~? Si vous observez un achat puis une vente, quelle est votre probabilité que la valeur soit 120~?
\end{question}

\begin{question}[$\star\star$ \iqroles{QT, QR} \iqfirm{teneur de marché, hedge fund} \iqlevel{junior, confirmé}]\label{iq:iv2:markets-and-microstructure:20}
Vous faites du market making, et votre achat moyen s'exécute à 10.00 alors que le prix milieu est de 10.01. Vous touchez un rabais de 0.2 cent. En moyenne, le prix milieu vaut 10.005 une seconde plus tard et 9.995 une minute plus tard. Quels sont vos profits par action, valorisés à ces deux instants~? Si vous traitez 1 million d'actions par jour, combien gagnez-vous à chaque horizon~? Auquel de ces chiffres accordez-vous du crédit, et que devez-vous faire~?
\end{question}

\begin{question}[$\star\star$ \iqroles{QT, Trader, Analyste} \iqfirm{hedge fund} \iqlevel{junior}]\label{iq:iv2:markets-and-microstructure:21}
Vous vendez à découvert 100 actions à 50~\$ pendant trois mois. Le prêteur facture des frais annuels de 6~\% sur la valeur, et vous percevez 4~\% par an sur le produit de la vente. Le titre verse un dividende de 0.50~\$ pendant cette période, et le cours reste inchangé. Quel est votre coût~? De combien le titre doit-il baisser pour que vous atteigniez le point d'équilibre~? Quelle est votre perte maximale possible~?
\end{question}

\begin{question}[$\star\star$ \iqroles{QT, Trader, Dev} \iqfirm{hedge fund, teneur de marché} \iqlevel{junior}]\label{iq:iv2:markets-and-microstructure:22}
Un titre américain a clôturé hier à 50~\$. Aujourd'hui il s'échange à 44.90~\$ et est coté avec un cours acheteur de 44.88~\$ et un cours vendeur de 44.90~\$. Que fait le test de prix de la vente à découvert au titre de la Rule 201~? Pendant combien de temps s'applique-t-il~? À quels prix pouvez-vous désormais vendre à découvert, et pouvez-vous toujours vendre au cours acheteur si vous détenez les actions~?
\end{question}

\begin{question}[$\star\star$ \iqroles{QT, Trader, Analyste} \iqfirm{hedge fund, banque} \iqlevel{junior}]\label{iq:iv2:markets-and-microstructure:23}
Une journée de négociation compte quatre périodes égales. Leurs parts du volume de la journée sont de 40~\%, 15~\%, 15~\% et 30~\%, et les prix moyens y sont de 100.0, 100.5, 101.0 et 100.2. Quel est le VWAP de la journée~? Vous achetez selon un échéancier pondéré par le temps, avec des montants égaux dans chaque période. Quel prix moyen obtenez-vous, et à quelle distance êtes-vous de l'indice de référence, en points de base~? Qu'aurait donné un échéancier qui suit le volume~?
\end{question}

\begin{question}[$\star\star$ \iqroles{QT, Trader, Analyste} \iqfirm{hedge fund, banque} \iqlevel{junior, confirmé}]\label{iq:iv2:markets-and-microstructure:24}
Votre ordre représente 30~\% du volume du titre aujourd'hui. Les autres ont échangé en moyenne à 100.00, et vous avez acheté en moyenne à 100.30, soit le coût de votre achat. Quel est le VWAP de la journée~? Que dit votre rapport par rapport à lui~? Quel est le coût réel, et que devient le rapport si vous représentez la totalité du volume~?
\end{question}

\begin{question}[$\star\star$ \iqroles{QT, Trader} \iqfirm{hedge fund, banque} \iqlevel{junior, confirmé}]\label{iq:iv2:markets-and-microstructure:25}
Vous devez vendre un gros bloc. Un dealer le rachète immédiatement à 0.8~\% sous le prix milieu. Si vous l'écoulez vous-même sur deux jours, le coût espéré est de 0.5~\%, avec un écart-type de 1.5~\%. Supposez que vous pénalisiez un choix de $\lambda$ fois sa variance, en unités de rendement. Pour quel $\lambda$ êtes-vous indifférent~? Que demanderiez-vous d'autre avant de décider~?
\end{question}

\begin{question}[$\star\star$ \iqroles{QT, QR, LLE, Dev} \iqfirm{teneur de marché, hedge fund} \iqlevel{junior, confirmé}]\label{iq:iv2:markets-and-microstructure:26}
Un cours vendeur périmé portant sur 100 actions reste affiché sur une bourse, à 2 cents sous la fair value. L'annulation d'un teneur de marché met un temps uniformément réparti entre 150 et 350 microsecondes pour arriver, et votre ordre met un temps uniformément réparti entre 200 et 400, indépendamment. Vous payez des frais de preneur de liquidité de 0.3~cent si vous exécutez, et vous ne payez rien si vous ratez. Quelle est la probabilité que vous l'emportiez, et que espérez-vous gagner par action et par tentative~?
\end{question}

\begin{question}[$\star\star$ \iqroles{QT, Trader} \iqfirm{hedge fund, banque} \iqlevel{junior, confirmé}]\label{iq:iv2:markets-and-microstructure:27}
Vous voulez acheter, et franchir la fourchette coûte 1 cent par rapport au prix milieu. Vous pouvez d'abord essayer un dark pool, qui exécute au prix milieu. S'il vous exécute, vous économisez le cent. Sinon, votre intention a fuité et le prix monte de 0.5~cent avant que vous ne franchissiez la fourchette. À partir de quelle probabilité d'exécution a-t-il intérêt à essayer d'abord~? Quelle est la valeur pour une probabilité de 40~\%~? Qu'est-ce qui manque à ce calcul~?
\end{question}

\begin{question}[$\star\star$ \iqroles{QT, Analyste} \iqfirm{teneur de marché, banque} \iqlevel{junior}]\label{iq:iv2:markets-and-microstructure:28}
Un dealer exécute les ordres de détail sur un titre dont la fourchette est de 2 cents. Il améliore chaque ordre de 0.2~cent sur le cours affiché, verse 0.1~cent par action au courtier, et, en moyenne, le prix évolue ensuite de 0.2~cent contre lui. Combien gagne-t-il par action~? Pourquoi peut-il payer pour ce flux tout en offrant mieux que le prix de la bourse~?
\end{question}

\begin{question}[$\star\star$ \iqroles{QR, QT, DS} \iqfirm{teneur de marché, hedge fund} \iqlevel{junior, confirmé}]\label{iq:iv2:markets-and-microstructure:29}
Un titre a une volatilité quotidienne réelle de 2~\% et une fourchette de 4 points de base. Une journée compte 390 minutes, et le prix des transactions rebondit entre le cours acheteur et le cours vendeur, indépendamment d'un échantillon à l'autre. Quelle volatilité quotidienne estimez-vous en sommant les rendements au carré à une minute~? Et à une seconde~? Que feriez-vous pour y remédier~?
\end{question}

\begin{question}[$\star\star$ \iqroles{QT, Trader, Risque} \iqfirm{hedge fund, banque} \iqlevel{junior}]\label{iq:iv2:markets-and-microstructure:30}
Vous achetez pour 10 millions de dollars d'un titre dont la volatilité annuelle est de 30~\%. Sa corrélation avec un indice de volatilité 15~\% est de 0.75. Quelle quantité de l'indice vendez-vous pour vous couvrir~? Quelle est la volatilité annuelle résiduelle, en pourcentage et en dollars, et quelle est-elle par jour, avec 16 comme racine carrée du nombre de jours de négociation dans l'année~?
\end{question}

\begin{question}[$\star\star$ \iqroles{QT, QR} \iqfirm{teneur de marché, hedge fund} \iqlevel{junior, confirmé}]\label{iq:iv2:markets-and-microstructure:31}
Si vous cotez $\delta$ cents par rapport au prix milieu, vous êtes exécuté à un taux de $e^{-k\delta}$ par seconde, avec $k = 0.5$ par cent. Ignorez la sélection adverse et l'inventaire. Quel $\delta$ maximise votre taux de profit espéré d'un seul côté~? Quels sont alors le taux d'exécution et le taux de profit~? Comment évoluent-ils si la concurrence double $k$~?
\end{question}

\begin{question}[$\star\star$ \iqroles{QT, Trader, QR} \iqfirm{hedge fund, teneur de marché} \iqlevel{junior, confirmé}]\label{iq:iv2:markets-and-microstructure:32}
Vous devez acheter maintenant ou dans la minute. La fourchette est de 2 cents, donc franchir celle-ci coûte 1 cent par rapport au prix milieu du jour. Si vous affichez plutôt au cours acheteur, vous êtes exécuté dans la minute avec une probabilité $p$. Sinon, le prix milieu aura en moyenne monté de 0.8~cent, et vous franchissez alors la fourchette. Mesuré par rapport au prix milieu du jour, quel montant espérez-vous payer~? Pour quelle valeur de $p$ a-t-il intérêt à afficher~? Faut-il aussi déduire quelque chose de la baisse du prix milieu au moment de l'exécution~?
\end{question}

\begin{question}[$\star\star$ \iqroles{QR, QT, DS} \iqfirm{teneur de marché, hedge fund} \iqlevel{junior, confirmé}]\label{iq:iv2:markets-and-microstructure:33}
Vous découpez le volume d'un titre en tranches de 100 actions chacune, et vous classez le volume d'achat et de vente de chaque tranche. Quatre tranches ont des achats et des ventes de 60 et 40, 70 et 30, 50 et 50, et 90 et 10. Quel est le VPIN, c'est-à-dire le déséquilibre absolu moyen des ordres par tranche, en proportion de la taille de la tranche~? Qu'est-il censé montrer~? Sur quoi faudrait-il être prudent~?
\end{question}

\begin{question}[$\star\star$ \iqroles{QT, Trader, Dev} \iqfirm{hedge fund, teneur de marché} \iqlevel{junior, confirmé}]\label{iq:iv2:markets-and-microstructure:34}
Un carnet affiche 100 actions à chaque niveau, à 1, 2, 3 cents et ainsi de suite au-dessus du prix milieu. Vous avez besoin de 1\,000 actions. Combien coûte un seul balayage au-dessus du prix milieu, en dollars~? Autrement, vous achetez 100 actions par minute pendant 10 minutes à un cent au-dessus du prix milieu, tandis que chaque tranche fait monter le prix milieu de 0.2~cent de façon durable. Combien cela coûte-t-il~? Si le titre varie de 5 cents par minute en écart-type, quel risque prennent les tranches~?
\end{question}

\section{Modèles soumis à questions}

\begin{question}[$\star\star\star$ \iqroles{QR, QT} \iqfirm{hedge fund, teneur de marché} \iqlevel{junior, confirmé}]\label{iq:iv2:markets-and-microstructure:35}
Vous vendez 1\,000 actions en deux tranches, $x_1$ puis $x_2$. L'impact temporaire propre à chaque tranche coûte $a x_i^2$. Votre première tranche fait baisser le prix auquel la seconde est exécutée, ce qui coûte $b x_1 x_2$. Détenir $x_2$ pendant toute la période ajoute une pénalité de risque $c x_2^2$. Prenez $a = 0.0002$, $b = 0.0001$ et $c = 0.0001$. Quelle répartition est la meilleure, et quel est le coût total~? Comparez-la avec une répartition égale et avec une exécution en une seule fois. Que se passe-t-il pour $c = 0$, et pour $c = 0$ et $b = 2a$~?
\end{question}

\begin{question}[$\star\star\star$ \iqroles{QR, QT, Trader} \iqfirm{teneur de marché, hedge fund} \iqlevel{junior, confirmé}]\label{iq:iv2:markets-and-microstructure:36}
Le meilleur cours vendeur porte 3 lots, et le meilleur cours acheteur en porte 7. Des ordres au marché d'achat et de vente d'un seul lot arrivent avec la même probabilité, et il n'arrive aucun nouvel ordre à cours limité. Le prix milieu monte d'un demi-tick si le cours vendeur est consommé en premier, et baisse d'un demi-tick si le cours acheteur l'est. Quelle est la probabilité qu'il monte~? Quel est le mouvement espéré~? Comparez-le au microprix, et expliquez pourquoi les deux diffèrent.
\end{question}

\begin{question}[$\star\star\star$ \iqroles{QR, QT} \iqfirm{teneur de marché, hedge fund} \iqlevel{junior, confirmé}]\label{iq:iv2:markets-and-microstructure:37}
Les signes des transactions, $+1$ pour un achat et $-1$ pour une vente, suivent une chaîne de Markov, de sorte que le signe suivant espéré vaut 0.8 fois le signe courant. Le prix ne bouge que sur les surprises, de 1 cent par unité de surprise. De combien le prix bouge-t-il sur un achat qui suit un achat~? Sur une vente qui suit un achat~? Quel est le mouvement total du prix après une série de cinq achats, en partant d'une position neutre~? Le flux d'ordres est fortement autocorrélé~: pourquoi les variations de prix ne le sont-elles pas~?
\end{question}

\begin{question}[$\star\star\star$ \iqroles{QT, Trader} \iqfirm{teneur de marché, banque} \iqlevel{junior, confirmé}]\label{iq:iv2:markets-and-microstructure:38}
Un client vous demande de lui racheter 1 million d'actions d'un seul coup. Le volume quotidien est de 2 millions, et la volatilité quotidienne de 2~\%. Vous les revendrez dans le temps à 20~\% du volume. Combien de temps cela prend-il~? Prenez votre coût moyen égal aux deux tiers de $\sigma\sqrt{Q/V}$. Quelle est sa valeur~? Quel est l'écart-type de votre résultat lorsque vous conservez une position qui se réduit~? Si vous facturez un demi-écart-type pour le risque, quelle décote cotez-vous~? Que fait une couverture par un indice, si la volatilité propre au titre est de 1.2~\% par jour~?
\end{question}

\begin{question}[$\star\star\star$ \iqroles{QR, QT} \iqfirm{hedge fund, teneur de marché} \iqlevel{confirmé}]\label{iq:iv2:markets-and-microstructure:39}
Vous détenez $q$ actions d'un titre au prix milieu $s$, avec une utilité exponentielle d'aversion au risque $\gamma$ et une variance du prix de $\sigma^2$ par unité de temps. Vous conserverez la position jusqu'à un instant $\tau$ à partir de maintenant. Déterminez le prix auquel vous êtes indifférent à l'achat d'une action supplémentaire, et de même pour la vente d'une action. Quelle est leur moyenne, le prix de réservation~? Avec $s = 100$, $q = 5$, $\gamma = 0.1$, $\sigma = 2$, $\tau = 0.5$ et une fourchette totale donnée par $\gamma\sigma^2\tau + (2/\gamma)\ln(1 + \gamma/k)$ avec $k = 1.5$, quelle cotation proposez-vous~?
\end{question}

\begin{question}[$\star\star\star$ \iqroles{QT, Trader, Dev, QD} \iqfirm{hedge fund, teneur de marché} \iqlevel{junior, confirmé}]\label{iq:iv2:markets-and-microstructure:40}
Un collègue a un algorithme qui place de gros ordres d'achat à quelques niveaux sous le marché, pour donner l'impression qu'il y a de la demande, et qui vend au cours vendeur une fois le prix monté. Il annule les ordres d'achat avant qu'ils ne soient exécutés. Un second algorithme annule 98~\% de ses ordres parce que ses cours ne cessent de changer. Est-ce la même chose~? Qu'est-ce qui, dans le droit américain, rend le premier illégal~? Comment le détecteriez-vous dans les ordres de votre propre société~?
\end{question}

\begin{question}[$\star\star\star$ \iqroles{QR, QT, Analyste} \iqfirm{teneur de marché, hedge fund} \iqlevel{confirmé}]\label{iq:iv2:markets-and-microstructure:41}
Un régulateur propose de réduire de moitié le pas de cotation minimal pour un titre valant 20~dollars, qui se négocie toujours avec une fourchette d'un tick et une file longue. Qui y gagne et qui y perd~: les investisseurs de long terme, les teneurs de marché proches de la tête de file, ceux proches de la fin de file, et les traders rapides qui passent devant~? Que mesureriez-vous ensuite pour l'évaluer~?
\end{question}

\begin{question}[$\star\star\star$ \iqroles{QR, QT, LLE} \iqfirm{teneur de marché, hedge fund} \iqlevel{junior, confirmé}]\label{iq:iv2:markets-and-microstructure:42}
Vous cotez une demi-fourchette de $h$ cents. Les traders ordinaires exécutent contre vous à un rythme d'une fois par seconde. Un signal public déplace la fair value de 2 cents à un rythme de 0.5 par seconde. Il y a 3 sociétés aussi rapides, l'une qui cote et deux qui s'emploient à exploiter la cotation~; chacune des trois a la même probabilité d'atteindre la bourse en premier. Quelle valeur de $h$ permet à la société qui cote d'être à l'équilibre~? Que devient $h$ avec une seule société, et avec un très grand nombre~? Serait-il utile de supprimer la course~?
\end{question}

\begin{question}[$\star\star\star$ \iqroles{QT, QR} \iqfirm{teneur de marché} \iqlevel{junior, confirmé}]\label{iq:iv2:markets-and-microstructure:43}
Un marché de contrats à terme exécute au prorata pur. D'autres affichent 1\,000 lots à votre prix. Une transaction porte sur 200 lots avec une probabilité de 90~\%, après quoi le prix ne bouge pas, ou sur 800 lots avec une probabilité de 10~\%, après quoi il bouge d'un tick contre vous. Vous gagnez un demi-tick par rapport à la fair value sur chaque exécution. Si vous affichez $q$ lots, combien espérez-vous gagner par transaction~? Que se passe-t-il lorsque $q$ augmente~? Si vous ne pouvez jamais exécuter plus de 200 lots en une seule transaction, combien affichez-vous~?
\end{question}

\begin{question}[$\star\star\star$ \iqroles{QR, DS, QT} \iqfirm{hedge fund, banque} \iqlevel{junior, confirmé}]\label{iq:iv2:markets-and-microstructure:44}
Un fonds répartit au hasard ses ordres entre deux courtiers, 400 pour chacun. Le coût moyen est de 3.0 points de base avec l'un et de 4.2 avec l'autre, et l'écart-type par ordre est de 30 points de base pour les deux. La différence est-elle significative~? Combien d'ordres par courtier faut-il pour détecter une différence de 1.2 point de base avec une puissance de 80~\% au seuil de 5~\%, en test bilatéral~? Quel rapport la taille du coût a-t-elle avec la manière dont il faut mener cette expérience~?
\end{question}

\begin{question}[$\star\star\star$ \iqroles{QR, QT} \iqfirm{hedge fund, teneur de marché} \iqlevel{junior, confirmé}]\label{iq:iv2:markets-and-microstructure:45}
Vous vendez 50\,000 actions dans chacune des deux périodes. Chaque action vendue fait baisser le prix de façon permanente de $b = \$0.000001$, et les transactions s'effectuent au prix qui suit les ventes de la période. Un prédateur qui connaît votre plan vend $y$ actions dans la première période en même temps que vous, et les rachète après votre deuxième période, lorsque ce rachat fait remonter le prix de $by$. Quelle valeur de $y$ maximise le profit du prédateur, et quel est ce profit~? Quel est votre surcoût, et quelle part représente-t-il de votre coût sans prédateur~? Qu'enseigne cela sur la manière dont vous devriez échelonner vos exécutions~?
\end{question}

\begin{question}[$\star\star\star$ \iqroles{QR, QT, Trader} \iqfirm{hedge fund, banque} \iqlevel{confirmé}]\label{iq:iv2:markets-and-microstructure:46}
Vous devez vendre un quart du volume d'une journée, avec une volatilité quotidienne de 2~\%. Supposez que, sur un horizon de $T$ jours, votre coût espéré soit $\tfrac23\sigma\sqrt{Q/(VT)}$ et que votre risque soit $\sigma\sqrt{T/3}$, et que vous minimisiez le coût plus une unité de risque. Quel horizon est le meilleur~? Quels sont les deux termes à cet horizon~? Que dit la loi empirique en racine carrée qui doit vous rendre prudent envers ce modèle de coût~?
\end{question}

\begin{question}[$\star\star\star$ \iqroles{QT, Trader, Risque} \iqfirm{teneur de marché, hedge fund} \iqlevel{junior, confirmé}]\label{iq:iv2:markets-and-microstructure:47}
Le cours de référence d'une action américaine est de 100~dollars, au milieu de la journée de négociation, et sa bande Limit Up--Limit Down est de 5~\%. À quels prix se situent les bandes~? Le titre se négocie jusqu'à la bande inférieure, et les cours vendeurs restent là, avec des cours acheteurs en dessous. Que se passe-t-il au bout de 15 secondes~? Comment la bande évolue-t-elle dans la première et la dernière partie de la journée~? Vous avez vendu un put sur ce titre et vous êtes couvert en delta~: quel est votre problème pendant la pause, et que faites-vous~?
\end{question}

\begin{question}[$\star\star\star\star$ \iqroles{QR} \iqfirm{hedge fund, teneur de marché} \iqlevel{confirmé}]\label{iq:iv2:markets-and-microstructure:48}
La valeur $v$ d'un actif suit une loi normale autour de $p_0$, de variance $\sigma_v^2$. Un seul initié observe $v$ et soumet un ordre $x$. Les noise traders soumettent $u$, une variable normale de moyenne nulle et de variance $\sigma_u^2$, indépendante de $v$. Un teneur de marché n'observe que $x + u$ et fixe le prix à $p_0 + \lambda(x + u)$, égal à l'espérance conditionnelle sachant le flux d'ordres. Déterminez la stratégie de l'initié $x = \beta(v - p_0)$ et $\lambda$. Avec $\sigma_v = 2$~dollars et $\sigma_u = 1\,000$ actions, quels sont $\beta$, $\lambda$, le profit espéré de l'initié et la variance résiduelle du prix~? Qui paie l'initié~?
\end{question}

\begin{question}[$\star\star\star\star$ \iqroles{QR} \iqfirm{hedge fund, banque, teneur de marché} \iqlevel{confirmé}]\label{iq:iv2:markets-and-microstructure:49}
Vous devez vendre $X$ actions avant l'instant $T$, en détenant $x(t)$ actions à l'instant $t$. Vendre à un rythme $v$ vous coûte $\eta v$ par action. Vous minimisez $\int_0^T \eta v^2\,dt + \lambda\sigma^2\int_0^T x^2\,dt$. Dérivez la trajectoire optimale. Soit $\kappa^2 = \lambda\sigma^2/\eta$. Quelle part est vendue à mi-parcours si $\kappa T = 1$, et si $\kappa T = 3$~? Dans chaque cas, comment le coût et la variance du prix se comparent-ils à une vente à rythme constant~? Quelle est la trajectoire pour $\lambda = 0$~?
\end{question}

\begin{question}[$\star\star\star\star$ \iqroles{QR} \iqfirm{hedge fund, teneur de marché} \iqlevel{confirmé}]\label{iq:iv2:markets-and-microstructure:50}
Vous vendez $X$ actions en trois transactions à des instants également espacés, de tailles $x_0$, $x_1$ et $x_2$. Une transaction de taille $x$ fait bouger le prix de $\kappa x$, ce qui vous coûte $\tfrac12\kappa x^2$ sur elle-même. Le mouvement antérieur se résorbe alors, et une fraction $a$ de celui-ci subsiste à l'instant de la transaction suivante. Le coût total est donc $\tfrac12\kappa\left(x_0^2 + x_1^2 + x_2^2\right) + \kappa a\left(x_0x_1 + x_1x_2\right) + \kappa a^2 x_0x_2$. Déterminez les tailles optimales. Avec $a = 0.5$, quelles sont les parts, et combien économisez-vous par rapport à des tiers égaux, et par rapport à une vente en une seule fois~? Que se passe-t-il lorsque $a$ tend vers 0 et vers 1~?
\end{question}
